% ID: MAT-GAN-PAR-001
% Titolo: Parabola con vertice assegnato e passante per un punto
% Disciplina: Matematica
% Area: Geometria analitica
% Argomento: Parabola
% Sottoargomento: Forma con vertice e asse parallelo all'asse y
% Classe: Terza liceo scientifico
% Difficolta: 3
% Tipo: problema_numerico
% Risultato: y = 2(x-1)^2 - 2; intersezioni con l'asse x in x = 0 e x = 2
% Tempo_stimato: 9 min
% Competenze: uso della forma canonica con vertice, calcolo del parametro, intersezioni con gli assi
% Prerequisiti: parabola, vertice, equazioni di secondo grado
% Tag: matematica, geometria_analitica, parabola, vertice, intersezioni
% Fonte: esempio_originale
% Autore: OpenAI Codex
% Licenza: CC-BY-SA-4.0

\begin{esercizio}
Determina l'equazione della parabola con asse parallelo all'asse \(y\),
vertice \(V(1,-2)\) e passante per il punto \(P(3,6)\).

Trova poi le intersezioni della parabola con l'asse \(x\).
\end{esercizio}

\begin{soluzione}
Una parabola con asse parallelo all'asse \(y\) e vertice \(V(x_V,y_V)\) ha
equazione:
\[
y=a(x-x_V)^2+y_V.
\]

Nel nostro caso:
\[
y=a(x-1)^2-2.
\]

Poiche la parabola passa per \(P(3,6)\), sostituiamo le coordinate del punto:
\[
6=a(3-1)^2-2.
\]

\[
6=4a-2.
\]

\[
4a=8,
\qquad a=2.
\]

L'equazione e quindi:
\[
y=2(x-1)^2-2.
\]

Per trovare le intersezioni con l'asse \(x\), imponiamo \(y=0\):
\[
0=2(x-1)^2-2.
\]

\[
(x-1)^2=1.
\]

\[
x-1=\pm 1.
\]

Quindi:
\[
x=0
\qquad \text{oppure} \qquad
x=2.
\]

Le intersezioni sono \((0,0)\) e \((2,0)\).
\end{soluzione}
